\chapter{Negative Results}
\label{ch:negative}

Several interventions improved an intermediate metric but harmed the complete
SLAM system. This chapter records those tests because they narrow the design
space and explain choices made in earlier chapters.

\section{Backbone adaptation}
\label{sec:neg:lora}

Encoder LoRA~\cite{hu2022lora} was tested as a broader alternative to
head-only adaptation. Reconstruction quality fell by $49\%$ and predicted
Gaussian scales increased by $1821\times$. Decoder-only LoRA improved by
$0.37$\dB{} after one epoch, then fell to $-1.61$\dB{} by epoch five while
the 99th-percentile scale increased by $12\times$.

Both failures combine a feature change with an exponentially parameterised
scale head. The measured signature motivates the frozen-backbone design in
\cref{ch:adapt}. It applies to the tested objectives and regularisation;
other joint adaptation methods remain open.

\section{Independent multi-pair averaging}
\label{sec:neg:aggregation}

One way to reduce pairwise scale variation is to predict a new keyframe
against several partners, align their depth scales, and average the maps.
With $m\in\{1,2,4,8\}$ predictions, scale disagreement falls from $5.16\%$
to $3.06\%$. A single sixteen-view VGGT forward reaches $1.68\%$.
Independent averaging reduces random variation but cannot impose one shared
scale across all views. It also requires $m$ forward passes per keyframe.

\section{Loop-closure retrieval refitting}
\label{sec:neg:retrieval}

The inherited retrieval system uses a PCA transform and an ASMK codebook.
A $2048$-centroid refit on $279$ TUM keyframes improves weighted offline
Recall@1 by $20\%$. End-to-end SLAM gives a different result:
\texttt{room} ATE improves by $8.2\%$, \texttt{360} degrades by $24.1\%$,
and \texttt{desk} degrades by $318.8\%$ to near the no-loop-closure control.
The evaluated system therefore retains the original MASt3R retrieval assets.
Candidate recall alone does not measure the effect of accepted edges on the
pose graph.

\section{Cross-family adaptation variation}
\label{sec:neg:attrib}

Head-only adaptation gains vary across datasets. Four explanations were
tested: residual mismatch, base-quality headroom, novelty to the training
corpus, and sensor modality. Each conflicts with at least one paired result
in \cref{tab:falsified}. The available evidence therefore reports the
variation without assigning one cause.

\section{Geometric scale correction}
\label{sec:neg:frontend}

Pairwise pointmaps show per-keyframe scale variation. We tested whether a
many-view model could estimate better local scales and whether those scales
would improve the Gaussian map.

\subsection{Context and backbone}
The gauge-free scale spread in \cref{tab:scale} separates model choice from
view context. The same VGGT~\cite{wang2025vggt} weights give $10.82\%$ spread
when run pairwise and $1.68\%$ when sixteen views enter one forward pass.
The two-view MASt3R control gives $5.16\%$. Multi-view context, rather than
the backbone alone, produces the lower spread.

\begin{table}[t]
\centering\small
\setlength{\tabcolsep}{5pt}
\begin{tabular}{@{}llrr@{}}
\toprule
Predictor & Context & Views & Disagreement \dn \\
\midrule
Splatt3R (deployed) & pairwise & 2 & $5.35\%$ \\
VGGT~\cite{wang2025vggt} & pairwise & 2 & $10.82\%$ \\
VGGT~\cite{wang2025vggt} & joint & 16 & \best{$1.68\%$} \\
\midrule
\multicolumn{4}{@{}l}{\emph{disjoint-window replication}} \\
Splatt3R & pairwise & 2 & $4.92\%$ \\
VGGT & joint & \phantom{0}8 & \best{$1.53\%$} \\
\bottomrule
\end{tabular}
\caption[Adjacent-view depth-scale disagreement]{Adjacent-view depth-scale disagreement, gauge-free, TUM freiburg1,
pooled over seven sequences ($84$ windows above the rule, $33$ disjoint
windows below it). Windows in the primary arm overlap by $80$--$97\%$, so the
result is replicated on eight-view windows tiled by exactly one span; the
replication is $7/7$ sequences, paired Wilcoxon $p=8.6\times10^{-8}$. VGGT's
$1.68\%$ sits at or below the RGB-D sensor's own view-dependent bias floor, so
it is a lower bound on how far coupling goes.}
\label{tab:scale}
\end{table}


\subsection{Map edit}
A sixteen-view estimate supplies one scale correction per keyframe cluster.
On TUM desk, sensor-referenced scale error improves from $6.48\%$ to
$1.02\%$. Rendering becomes worse: PSNR falls from $12.35$ to
$11.50$\dB{} and LPIPS rises from $0.507$ to $0.523$.
After $300$\,s refinement, PSNR changes from $13.81$ to $13.08$\dB{}.
The same trend appears on \texttt{360} and \texttt{room}.

\subsection{Scale prior during refinement}
Using the correction as a prior instead of a fixed edit also reduces image
quality. As prior weight rises from $0.05$ to $0.5$, PSNR falls from
$13.74$ to $13.38$\dB{} and LPIPS rises from $0.449$ to $0.456$.
The fitted scale spread remains $2.6\%$, compared with the reference
correction of $8.0\%$.

These tests show a conflict between isolated scale accuracy and the current
photometric map fit. Camera poses, local geometry, and appearance have been
estimated together; editing one part can reduce their internal agreement.
A prior explanation based on a $0.057$\,m keyframe baseline was withdrawn
after contribution-weighted rendering rays showed a baseline near $0.78$\,m.
The negative result rests on the measured edits and priors, not on that
earlier explanation.

\section{Result of the negative studies}
\label{sec:neg:method}

The failed routes improve reconstruction loss, retrieval recall, or local
scale accuracy without improving the final map or trajectory. For this
system, an intervention is retained only when its end-to-end metric improves
under the target execution path. This rule leads to three concrete choices:
freeze shared features, keep the inherited retrieval assets, and avoid
map-only scale correction.
